% !TEX program = lualatex
\documentclass[12pt]{article}

% LuaLaTeX: fontspec + polyglossia для корректной работы с кириллицей
\usepackage{fontspec}
\defaultfontfeatures{Ligatures=TeX,Scale=MatchLowercase}
\setmainfont{Times New Roman}
\newfontfamily\cyrillicfont{Times New Roman}
\usepackage{polyglossia}
\setmainlanguage{russian}

\usepackage[style=numeric, backend=biber]{biblatex} % Обязательный пакет для корректной сборки
\usepackage{graphicx}        						% Графика
\usepackage{microtype}       						% Микротипографика
\usepackage{amssymb}        						% Математические символы
\usepackage{amsmath}        						% Математика
\usepackage{tikz}           						% Рисование
\usepackage{circuitikz}     						% Электрические схемы
\usepackage{fix-cm}         						% Чтобы работал большой размер шрифта
\usepackage{anyfontsize}    						% Для любого размера шрифта
\usepackage{ragged2e}       						% Для команды \justifying

\usepackage[a4paper,left=20mm,right=20mm,top=20mm,bottom=20mm]{geometry}

\newfontfamily\titletimesnewroman[Scale=1.4]{Times New Roman}

\begin{document}

\begin{center}
	
Федеральное государственное бюджетное образовательное учреждение высшего\\
профессионального образования\\
«Национальный исследовательский институт МИЭТ»
\vskip2cm
кафедра электротехники
\vskip2cm
% локально убрать отступ первой строки (parindent=0) внутри minipage, шрифт сохраняется
{
  \setlength{\parindent}{0pt}
  \hbadness=10000 % отключить предупреждения underfull
  \linespread{1.5}\selectfont % увеличить межстрочный интервал
  % первая строка — центр
  \centering\titletimesnewroman\bfseries
  Пояснительная записка к курсовому проекту\par
  % следующие строки — растянуть на всю ширину страницы
  \titletimesnewroman\bfseries
  \makebox[\textwidth][s]{на тему: «Линейные электрические цепи постоянного и}\par
  \makebox[\textwidth][s]{переменного тока. (Методы расчета). Переходные процессы.»}\par
}

\end{center}
\end{document}